\chapter{Positional and Keyword Arguments}

\section{Introduction}

\begin{tcolorbox}[title=Goal]
\textbf{Refactor} the Session 4 classifier so the same pipeline can be controlled from function calls. In this session, students learn how positional arguments depend on order, how keyword arguments improve readability, and how default values let one function support multiple run modes.
\end{tcolorbox}

\begin{tcolorbox}[title=You will work in]
Use the following files for this session:
\begin{itemize}
    \item \texttt{session5/session5\_iris\_template.py} (student template)
    \item \texttt{session5/iris\_data.csv} (dataset)
    \item \texttt{session5/helper/iris\_utils.py} (provided helper module)
\end{itemize}
\textbf{How to run the file:} Use the terminal command: \\
\texttt{python3 session5/session5\_iris\_template.py} \\
\textbf{Do not edit the helper module.} We will only import and use it.
\end{tcolorbox}

In the previous sessions (session 4), you wrote several \textbf{functions} and called them from a \texttt{main()} function.
That already gave you a powerful skill. It was your the first major refactor: instead of keeping everything in one long script, students learned to divide the work into smaller parts such as prediction, counter updates, and summary printing.

In this Session 5, we continues that same refactoring mindset (i.e., divide the work into smaller parts).but we make your functions \textbf{more flexible} by controlling them through \textbf{arguments} when you call them.

\begin{tcolorbox}[title=Why this matters (Real-World Hook)]
Parameterizing functions is akin to how APIs expose configuration options (e.g., passing options to a machine-learning model). It allows users to control the behavior of your code without having to rewrite its internal logic.
\end{tcolorbox}

This is an important programming idea. Good functions do not only hide logic. They also expose clean ways for a caller to control that logic.
Because the positional section is for reading only, the implementation tasks begin in the keyword-argument section.

\section{What You Will Learn}
\textbf{In this session, you will learn to:}
\begin{enumerate}\setlength{\itemsep}{5pt}
\item \textbf{Use positional arguments carefully}: understand how Python matches values by order.
\item \textbf{Use keyword arguments clearly}: understand how Python matches values by parameter name.
\item \textbf{Add default values}: make functions easier to call while keeping sensible behavior.
\item \textbf{Build a small API-style pipeline}: let one setup function control loading data, applying thresholds, printing each sample, and using a fitted threshold.
\item \textbf{Debug argument flow line by line}: inspect how values move from \texttt{main()} into helper functions.
\end{enumerate}

\subsection{Parameters and Arguments}
These two words are often confused by beginners, so we state the difference clearly:
\begin{itemize}
    \item A \textbf{parameter} is the name written in the function definition.
    \item An \textbf{argument} is the real value passed into that parameter when the function is called.
\end{itemize}

\begin{minted}[fontsize=\small, linenos, breaklines=true, frame=lines]{python}
def make_print_status(status_text):
    print(f"[STATUS] {status_text}")

make_print_status("Build dataset")
\end{minted}

In the example above:
\begin{itemize}
    \item \texttt{status\_text} is the parameter,
    \item \texttt{"Build dataset"} is the argument.
\end{itemize}

\subsection{Why Session 5 Matters}
In a real project, it is common to reuse the same function many times with slightly different settings. If a function only works for one hardcoded situation, it becomes difficult to test, difficult to reuse, and difficult to extend.

Session 5 therefore teaches students a second level of refactoring:
\begin{itemize}
    \item Session 4 refactored \textbf{structure} by moving code into functions.
    \item Session 5 refactors \textbf{interfaces} by making those functions easier to call in different ways.
\end{itemize}

\section{Session 4 to Session 5 Refactor Map}
The classifier concept remains the same, but the surrounding structure becomes more flexible:
\begin{itemize}
    \item Session 4 static dataset setup $\rightarrow$ Session 5 CSV loading with \texttt{load\_iris\_data(filepath=None)}
    \item Session 4 fixed prediction loop $\rightarrow$ Session 5 configurable loop with \texttt{initialize\_predictions(dataset, settings, print\_each=True)}
    \item Session 4 positional accuracy calculation $\rightarrow$ Session 5 keyword-friendly \texttt{calculate\_accuracy(correct=0, total=0)}
    \item Session 4 direct top-level flow $\rightarrow$ Session 5 orchestration through \texttt{setup\_application(filepath=None, threshold=None, print\_each=True, use\_fit=False)}
    \item Session 4 one program run $\rightarrow$ Session 5 multiple program runs from the same \texttt{main()}
\end{itemize}

\section{Canonical Function Set}
Complete these functions in \texttt{session5\_iris\_template.py}. The table below shows the Session 5 target structure:

\begin{landscape}
\small
\renewcommand{\arraystretch}{0.3}

\begin{longtable}{|p{0.3\linewidth}|p{0.4\linewidth}|p{0.2\linewidth}|}
\caption{Session 5 canonical functions: what to implement, what stays the same from Session 4, and where positional, keyword, and mixed calls are introduced.} \\
\hline
\textbf{Section / Function} & \textbf{Code to put inside the function} & \textbf{Session 5 note / modification} \\
\hline
\endfirsthead

\multicolumn{3}{c}%
{{\bfseries \tablename\ \thetable{} -- continued from previous page}} \\
\hline
\textbf{Section / Function} & \textbf{Code to put inside the function} & \textbf{Session 5 note / modification} \\
\hline
\endhead

\hline \multicolumn{3}{|r|}{{Continued on next page}} \\ \hline
\endfoot

\hline \hline
\endlastfoot

\texttt{make\_print\_status(status\_text)} & \begin{minted}[fontsize=\small, linenos, breaklines, xleftmargin=2em]{python}
print(f"[STATUS] {status_text}")
\end{minted}
& Same as Session 4. \\
\hline
\texttt{setup\_application\_list()} & \begin{minted}[fontsize=\small, linenos, breaklines, xleftmargin=2em]{python}
return [
    {"id": "flower1", "petal_length": 1.4, "species": "setosa"},
    ...
]
\end{minted}
& Same role as Session 4 and Session 6: provide the dataset for the lesson. \\
\hline
\texttt{compute\_threshold\_prediction(sample, threshold)} & \begin{minted}[fontsize=\small, linenos, breaklines, xleftmargin=2em]{python}
if sample["petal_length"] < threshold:
    return POSITIVE_LABEL
return NEGATIVE_LABEL
\end{minted}
& Same prediction rule as Session 4 and Session 6. \\
\hline
\texttt{derive\_true\_label(sample)} & \begin{minted}[fontsize=\small, linenos, breaklines, xleftmargin=2em]{python}
if sample["species"] == POSITIVE_LABEL:
    return POSITIVE_LABEL
return NEGATIVE_LABEL
\end{minted}
& Same idea as Session 4 and Session 6. \\
\hline
\texttt{update\_result\_counts(correct, wrong, total, y\_pred\_list, y\_pred, y\_true)} & \begin{minted}[fontsize=\small, linenos, breaklines, xleftmargin=2em]{python}
if y_pred == y_true:
    correct += 1
else:
    wrong += 1
total += 1
y_pred_list.append(y_pred)
return correct, wrong, total, y_pred_list
\end{minted}
& Same counter logic as Session 4 and Session 6. \\
\hline
\texttt{run\_prediction\_loop(dataset, threshold=2.0, print\_each=True)} & \begin{minted}[fontsize=\small, linenos, breaklines, xleftmargin=2em]{python}
for sample in dataset:
    y_pred = compute_threshold_prediction(sample, threshold)
    y_true = derive_true_label(sample)
    correct, wrong, total, y_pred_list = update_result_counts(
        correct, wrong, total, y_pred_list, y_pred=y_pred, y_true=y_true
    )
    if print_each:
        print(...)
\end{minted}
& Same loop structure as Session 4 and Session 6, but now used to teach mixed positional+keyword calls. \\
\hline
\texttt{calculate\_accuracy(correct=0, total=0)} & \begin{minted}[fontsize=\small, linenos, breaklines, xleftmargin=2em]{python}
if total > 0:
    accuracy = (correct / total) * 100
else:
    accuracy = 0.0
return accuracy
\end{minted}
& \textbf{Refactor in S5:} same accuracy logic, but now with keyword-friendly defaults. \\
\hline
\texttt{run\_classifier\_pipeline(threshold=2.0, print\_each=True)} & \begin{minted}[fontsize=\small, linenos, breaklines, xleftmargin=2em]{python}
dataset = setup_application_list()
correct, wrong, total, y_pred_list = run_prediction_loop(
    dataset, threshold=threshold, print_each=print_each
)
accuracy = calculate_accuracy(correct=correct, total=total)
return correct, wrong, total, y_pred_list, accuracy
\end{minted}
& \textbf{New in S5:} renamed wrapper with a clearer purpose and a mixed-call example. \\
\hline
\texttt{print\_summary(correct, wrong, total, y\_pred\_list, accuracy)} & \begin{minted}[fontsize=\small, linenos, breaklines, xleftmargin=2em]{python}
print("Correct:", correct)
print("Wrong:", wrong)
print("Total:", total)
print("Accuracy (%):", round(accuracy, 2))
print("All predictions:", y_pred_list)
\end{minted}
& Same summary idea as Session 4 and Session 6. \\
\hline
\texttt{main()} & \begin{minted}[fontsize=\small, linenos, breaklines, xleftmargin=2em]{python}
run default call
run positional override call
run keyword override call
run reordered keyword call
\end{minted}
& \textbf{New in S5:} demonstrates the same pipeline in several calling styles. \\
\hline
\end{longtable}
\end{landscape}


\begin{tcolorbox}[title=Good Practice]
This session is intentionally verbose. When students are learning argument passing, confusion usually comes from two places: not knowing \textbf{which function header changed}, and not knowing \textbf{which call site must also change}. The instructions below therefore say both where to edit and what code to put there.
\end{tcolorbox}

\section{Working with Keyword Arguments}
Keyword arguments solve a common readability problem. When a function accepts many parameters, a long positional call can become difficult to read. Unlike positional arguments, keyword arguments allow you to specify exactly which parameter is being assigned a value, enhancing readability and reducing errors. A keyword call makes the meaning visible because the parameter names appear directly in the function call.

In this section, students will keep the same classifier behavior but improve the way the program is controlled from outside the function. This is where the implementation tasks begin.


\subsubsection{Advantages of Keyword Arguments:}

There are several advantages:
\begin{itemize}\setlength{\itemsep}{0pt}
    \item Enhanced readability and clarity in function calls.
    \item Flexible argument ordering, reducing potential errors.
    \item Easy handling of optional parameters with default values.

\subsubsection{Keyword Argument Syntax}
When you use keyword arguments, you specify the name of the parameter followed by an equal sign and the value it should receive within the function call. For example,
\begin{center}
    width=1000
\end{center}

\noindent
The benefit of this approach is that you can rearrange the order of arguments when calling the function without affecting its behavior, as long as all required arguments are provided.

So far, we have not been using keyword arguments in our code, but we will start introducing them in this session to give you more control over the behavior of your functions without having to change their internal logic.


\subsection{Task 1: Refactor \texttt{calculate\_accuracy} to use keyword-friendly defaults}

At first, this may seem unnecessary because \texttt{calculate\_accuracy} only uses two parameters. However, this is still a very useful practice for beginners. The goal of this task is not only to make this one function work, but also to help you understand how keyword arguments behave in Python.

When a function uses keyword-friendly parameters:
\begin{itemize}
    \item the function call becomes easier to read,
    \item the meaning of each value becomes clearer,
    \item and the function becomes safer to use later when the program grows larger.
\end{itemize}

In this task, you will make a small refactor:
\begin{itemize}
    \item you will update the function header so both parameters have default values,
    \item and then you will call the function using keyword arguments instead of a positional-style call.
\end{itemize}

Do not change the main accuracy logic. The purpose here is to improve the function interface, not to redesign how the calculation works.

\begin{tcolorbox}[title=Task 1: Refactor \texttt{calculate\_accuracy} to use keyword-friendly defaults]
Work only in \texttt{session5/session5\_iris\_template.py}. Do not create a new file, and do not rename the function.

\begin{enumerate}
    \item Find the function named \texttt{calculate\_accuracy}. Read the function carefully before editing anything so you understand what it already does.
    \item Locate the \texttt{<your\_code>} marker in the function header.
    \item Replace only the function header part so it becomes:
    \texttt{def calculate\_accuracy(correct=0, total=0):}
    \item Notice what this change means:
    \begin{itemize}
        \item \texttt{correct=0} gives the parameter \texttt{correct} a default value of \texttt{0},
        \item \texttt{total=0} gives the parameter \texttt{total} a default value of \texttt{0},
        \item so Python can still call the function even if no argument is passed in.
    \end{itemize}
    \item Keep the existing calculation logic and return statement inside the function body. Do not remove the accuracy logic that is already there.
    \item After updating the function definition, go to \texttt{setup\_application(...)}.
    \item Find the place where \texttt{calculate\_accuracy(...)} is called.
    \item Change that call so it uses keyword arguments:
    \texttt{correct=metrics["correct"]} and \texttt{total=metrics["total"]}.
    \item Read the new function call slowly:
    \begin{itemize}
        \item the value from \texttt{metrics["correct"]} is being assigned to the parameter named \texttt{correct},
        \item the value from \texttt{metrics["total"]} is being assigned to the parameter named \texttt{total}.
    \end{itemize}
    \item Quick check: \texttt{calculate\_accuracy(correct=8, total=10)} should return \texttt{80.0}.
    \item Quick check: \texttt{calculate\_accuracy()} should not crash. Because the function now has default values, Python can call it even when no arguments are given.
\end{enumerate}
\end{tcolorbox}

Use this exact function shape as the target:

\begin{minted}[fontsize=\small, linenos, breaklines=true, frame=lines]{python}
def calculate_accuracy(correct=0, total=0):
    # <your_code>: keep the accuracy logic and return value here
\end{minted}

\begin{tcolorbox}[title=First Unit Test for Task 1]
If this is your first time using \texttt{unittest}, think of it as an automatic checker for your Python code. Instead of checking your function by hand every time, Python can run several checks for you and report whether your function behaves correctly.

For Task 1, use the dedicated test file:
\texttt{session5/test\_calculate\_accuracy.py}

This test file checks only \texttt{calculate\_accuracy}. That is useful because you may still be working on the rest of \texttt{session5/session5\_iris\_template.py}.

Run the test from the project root folder with this command:

\begin{minted}[fontsize=\small, linenos, breaklines=true, frame=lines]{bash}
python3 -m unittest session5.test_calculate_accuracy -v
\end{minted}

What each part means:
\begin{itemize}
    \item \texttt{python3} starts Python,
    \item \texttt{-m unittest} runs Python's built-in testing tool,
    \item \texttt{session5.test\_calculate\_accuracy} points to the test file,
    \item \texttt{-v} means ``verbose'' and shows each test name clearly.
\end{itemize}

How to read the result:
\begin{itemize}
    \item \texttt{OK} means all tests passed,
    \item \texttt{FAIL} means the function ran, but at least one answer was wrong,
    \item \texttt{ERROR} usually means Python could not read or run the function yet.
\end{itemize}

For this task, the most important checks are:
\begin{itemize}
    \item \texttt{calculate\_accuracy(correct=8, total=10)} should return \texttt{80.0},
    \item \texttt{calculate\_accuracy()} should return \texttt{0.0} without crashing,
    \item keyword calls and positional calls should match when the numbers are the same.
\end{itemize}
\end{tcolorbox}

This code means:
\begin{itemize}
    \item the function is still named \texttt{calculate\_accuracy},
    \item the parameter \texttt{correct} defaults to \texttt{0},
    \item the parameter \texttt{total} defaults to \texttt{0},
    \item and the existing logic inside the function should stay responsible for computing and returning the result.
\end{itemize}

Then, inside \texttt{setup\_application(...)}, update the function call to use keyword arguments:

\begin{minted}[fontsize=\small, linenos, breaklines=true, frame=lines]{python}
accuracy = calculate_accuracy(
    correct=metrics["correct"],
    total=metrics["total"],
)
\end{minted}

Study this call carefully:
\begin{itemize}
    \item \texttt{correct=} tells Python exactly which parameter should receive \texttt{metrics["correct"]},
    \item \texttt{total=} tells Python exactly which parameter should receive \texttt{metrics["total"]},
    \item so the call is easier to read than simply writing two bare values in order.
\end{itemize}

The main advantage of this approach is that the parameter names are visible at the call site. A reader does not need to guess what each value means. Another advantage is that, when using keyword arguments, the order can be changed without changing the behavior, as long as the correct parameter names are used.

For example, you can call \texttt{calculate\_accuracy(total=total, correct=correct)} and it will work just fine.




\begin{tcolorbox}[title=Line-by-Line Debugging for Task 1]
Set a breakpoint on the line that calls \texttt{calculate\_accuracy(...)}.
\begin{itemize}
    \item Inspect \texttt{metrics["correct"]} and \texttt{metrics["total"]}.
    \item Use \textbf{Step Into} and confirm that those values arrive in the parameters \texttt{correct} and \texttt{total}.
    \item Tell students to notice that keyword argument order can change without changing the result.
\end{itemize}
This is one of the simplest examples for teaching the difference between positional matching and keyword matching.
\end{tcolorbox}




\subsection{Task 2: Build the keyword-friendly \texttt{setup\_application} function}
It is also possible to have many keyword arguments in a function.
For example, in the \texttt{setup\_application} function, you can have multiple keyword arguments to control different aspects of ... 

Recall that in previous session the setup application function accept 6 positional arguments, which means the order of the arguments matters. 

\begin{tcolorbox}[title=Task 2: Build the keyword-friendly \texttt{setup\_application} function]
In the template, find \texttt{setup\_application}.

Replace the \texttt{<your\_code>} marker in the function header so all four parameters become optional and keyword-friendly:
\texttt{def setup\_application(filepath=None, threshold=None, print\_each=True, use\_fit=False):}
\end{tcolorbox}

Use this exact function shape as the target:

\begin{minted}[fontsize=\small, linenos, breaklines=true, frame=lines]{python}
def setup_application(filepath=None, threshold=None, print_each=True, use_fit=False):
    # <your_code>: load the dataset, read the default settings, and apply the threshold rules
\end{minted}

Inside the function, the threshold logic must follow this priority order:
\begin{enumerate}
    \item if \texttt{use\_fit} is \texttt{True}, learn the threshold from \texttt{fit\_threshold\_from\_setosa(...)},
    \item else if \texttt{threshold} is not \texttt{None}, use the manual threshold,
    \item else keep the default threshold from \texttt{get\_default\_settings()}.
\end{enumerate}

If \texttt{filepath} is omitted, use \texttt{session5/iris\_data.csv} so the default run works without extra arguments.

Put this logic inside the function:

\begin{minted}[fontsize=\small, linenos, breaklines=true, frame=lines]{python}
dataset = load_iris_data(filepath)
settings = get_default_settings()

# <your_code>: apply the threshold priority order here

metrics = initialize_predictions(dataset, settings, print_each)
accuracy = calculate_accuracy(
    correct=metrics["correct"],
    total=metrics["total"],
)
return settings, dataset, metrics, accuracy
\end{minted}

This function is the main refactoring target of Session 5. It turns several separate steps into one reusable entry point.

\begin{tcolorbox}[title=Checkpoint]
After Task 2, students should be able to explain the difference between these two cases:
\begin{itemize}
    \item \texttt{threshold=2.2} means ``use this manual threshold'',
    \item \texttt{use\_fit=True} means ``ignore the manual override and learn the threshold from the dataset''.
\end{itemize}
\end{tcolorbox}

\begin{tcolorbox}[title=Line-by-Line Debugging for Task 2]
This is the most important debugging task in Session 5.
\begin{itemize}
    \item Set a breakpoint inside \texttt{setup\_application(...)} before the threshold logic.
    \item Inspect \texttt{threshold}, \texttt{use\_fit}, and \texttt{settings["threshold"]}.
    \item Step line by line through the \texttt{if use\_fit} and \texttt{elif threshold is not None} branches.
    \item Ask students: ``Which branch ran, and why?''
\end{itemize}
This question forces students to connect the function call in \texttt{main()} with the branch behavior inside the function.
\end{tcolorbox}

\subsection{Task 3: Add the default run in \texttt{main()}}
\begin{tcolorbox}[title=Task 3: Add the default run in \texttt{main()}]
In the template, find the default-run block and replace the \texttt{<your\_code>} marker by uncommenting the call.

The run should use the simplest possible call:
\texttt{setup\_application()}
\end{tcolorbox}

\begin{minted}[fontsize=\small, linenos, breaklines=true, frame=lines]{python}
# <your_code>: uncomment the default run block below
# make_print_status("Default run")
# settings, dataset, metrics, accuracy = setup_application()
# print_summary("Default run", settings, metrics, accuracy)
\end{minted}

This run is important because it proves that the defaults are working together:
\begin{itemize}
    \item default filepath,
    \item default threshold,
    \item default \texttt{print\_each=True},
    \item default \texttt{use\_fit=False}.
\end{itemize}

\begin{tcolorbox}[title=Checkpoint]
If the default run fails, students should first check the function headers. In beginner code, many default-run bugs come from forgetting to add \texttt{=None}, \texttt{=True}, or \texttt{=False} in the function definition.
\end{tcolorbox}

\subsection{Task 4: Add the positional override run in \texttt{main()}}
\begin{tcolorbox}[title=Task 4: Add the positional override run in \texttt{main()}]
In the template, find the positional run block and replace the \texttt{<your\_code>} placeholder by uncommenting the call.

Students should use a positional call like this:
\end{tcolorbox}

\begin{minted}[fontsize=\small, linenos, breaklines=true, frame=lines]{python}
# <your_code>: uncomment the positional override run below
# make_print_status("Positional override run")
# settings, dataset, metrics, accuracy = setup_application(
#     "session5/iris_data.csv", 1.8, False, False
# )
# print_summary("Positional override run", settings, metrics, accuracy)
\end{minted}

Explain to students that the order of positional arguments matters:
\begin{itemize}
    \item first argument $\rightarrow$ \texttt{filepath}
    \item second argument $\rightarrow$ \texttt{threshold}
    \item third argument $\rightarrow$ \texttt{print\_each}
    \item fourth argument $\rightarrow$ \texttt{use\_fit}
\end{itemize}

\begin{tcolorbox}[title=Checkpoint]
If students accidentally swap \texttt{1.8} and \texttt{False}, the run will still be syntactically valid, but the behavior will be wrong. This is exactly why positional calls become harder to read as the number of arguments grows.
\end{tcolorbox}

\begin{tcolorbox}[title=Line-by-Line Debugging for Task 4]
Set a breakpoint on the positional call in \texttt{main()} and use \textbf{Step Into}.
\begin{itemize}
    \item Before stepping in, read the four argument values left to right.
    \item After stepping in, compare them with the parameter names in \texttt{setup\_application(...)}.
    \item Inspect \texttt{filepath}, \texttt{threshold}, \texttt{print\_each}, and \texttt{use\_fit}.
\end{itemize}
This side-by-side comparison is one of the best debugging exercises in the whole session.
\end{tcolorbox}

\subsection{Task 5: Add the keyword override run in \texttt{main()}}
\begin{tcolorbox}[title=Task 5: Add the keyword override run in \texttt{main()}]
In the template, find the keyword-run block and replace the \texttt{<your\_code>} marker by uncommenting the call.

Students should use a named call like this:
\end{tcolorbox}

\begin{minted}[fontsize=\small, linenos, breaklines=true, frame=lines]{python}
# <your_code>: uncomment the keyword override run below
# make_print_status("Keyword override run")
# settings, dataset, metrics, accuracy = setup_application(
#     threshold=2.2,
#     print_each=False,
# )
# print_summary("Keyword override run", settings, metrics, accuracy)
\end{minted}

Explain to students that the order of keyword arguments does not matter as long as the parameter names are correct. For example, the following is also valid:

\begin{minted}[fontsize=\small, linenos, breaklines=true, frame=lines]{python}
setup_application(
    print_each=False,
    threshold=2.2,
)
\end{minted}

\begin{tcolorbox}[title=Checkpoint]
Students should notice that the keyword override run is easier to read than the positional override run. The call itself explains the meaning of each changed value.
\end{tcolorbox}

\subsection{Task 6: Add the fit-based run in \texttt{main()}}
\begin{tcolorbox}[title=Task 6: Add the fit-based run in \texttt{main()}]
In the template, find the fit-based block and replace the \texttt{<your\_code>} marker by uncommenting the call.

Students should use:
\end{tcolorbox}

\begin{minted}[fontsize=\small, linenos, breaklines=true, frame=lines]{python}
# <your_code>: uncomment the fit-based run below
# make_print_status("Fit-based run")
# settings, dataset, metrics, accuracy = setup_application(
#     print_each=False,
#     use_fit=True,
# )
# print_summary("Fit-based run", settings, metrics, accuracy)
\end{minted}

This run introduces a very important programming idea: one flag can change the path of the program. The data loading stays the same. The prediction loop stays the same. Only the threshold selection behavior changes.

\begin{tcolorbox}[title=Checkpoint]
After Task 6, students should compare:
\begin{itemize}
    \item the default run threshold,
    \item the positional override threshold,
    \item the keyword override threshold,
    \item and the fit-based threshold.
\end{itemize}
This comparison helps students see that Session 5 is about controlling one pipeline in multiple ways.
\end{tcolorbox}

\begin{tcolorbox}[title=Common Argument Mistakes]
Students often hit the same beginner errors in this session:
\begin{itemize}
    \item \textbf{Missing required positional argument}: the function header was not updated with a default value.
    \item \textbf{SyntaxError: positional argument follows keyword argument}: a positional value was placed after a keyword argument in the same call.
    \item \textbf{TypeError: multiple values for argument}: the same parameter was filled once positionally and once by keyword.
    \item \textbf{Wrong branch in \texttt{setup\_application(...)}:} students expected \texttt{threshold=} to run, but \texttt{use\_fit=True} took priority.
\end{itemize}
\end{tcolorbox}

\section{Positional Arguments}
Positional arguments are the most common way to pass values to functions. You've already been using them without explicitly realizing it since session 3 of our tutorial session.
 
In positional arguments, the order of the arguments provided in the function call directly corresponds to the order defined in the function signature.

\subsection{Loading Iris Data with a Single Positional Argument}

Take for example when we are loading data from a file using the \texttt{load\_iris\_data} function. When we pass the file path \texttt{('..session4/\_iris\_data.csv')} as an argument to the function, we're utilizing what's known as a positional argument.

\begin{minted}[fontsize=\small, linenos, breaklines=true]{python}
def setup_application(filepath):
    """Load, configure, and run prediction pipeline."""
    dataset = load_iris_data(filepath)

\end{minted}

The corresponding function definition is:

\begin{minted}[fontsize=\small, linenos, breaklines=true]{python}
def load_iris_data(filepath):
    dataset = []
    with open(filepath, "r", encoding="utf-8", newline="") as csv_file:
        reader = csv.DictReader(csv_file)
        for row in reader:
            dataset.append(
                {
                    "id": row["id"],
                    "sepal_length": float(row["sepal_length"]),
                    "sepal_width": float(row["sepal_width"]),
                    "petal_length": float(row["petal_length"]),
                    "petal_width": float(row["petal_width"]),
                    "species": row["species"],
                }
            )
\end{minted}


Here, the parameter \texttt{filepath} is a positional argument, which means the order matters. The first argument passed during the function call will always be matched with the first parameter defined in the function signature.


\subsection{Initializing initialize predictions with Two Positional Arguments}

Another example is the function \texttt{initialize\_prediction}, defined with two positional parameters: \texttt{dataset} and \texttt{setting}.
When invoking this function, you must ensure that the arguments are provided in the correct sequence. Swapping these arguments inadvertently will cause unexpected behavior or errors in your program.

Thus, clear parameter naming and detailed documentation are essential to prevent such issues.



Calling \texttt{initialize\_robots(canvas, data)} with the \texttt{canvas} and \texttt{data} in the correct order is crucial. Swapping these arguments would lead to errors. Hence, careful naming and detail doct string can help prevent such mistakes.

\subsection{ Update metrics Using Multiple Positional Arguments}
% update_metrics(metrics, y_pred, y_true)

Let's examine a function that requires multiple positional arguments, \texttt{update\_metrics}, which is responsible for rendering robots on the canvas:


\begin{minted}[fontsize=\small, linenos, breaklines=true]{python}
def update_metrics(metrics, y_pred, y_true):
    """Update correct/wrong/total counters and append predictions.
    
    Args:
        metrics (dict): The metrics dictionary tracking accuracy.
        y_pred (str): The predicted label.
        y_true (str): The true label.
        
    Side Effects:
        Modifies the `metrics` dictionary in place.
    """
    is_correct = evaluate_prediction(y_pred, y_true)
    if is_correct:
        metrics["correct"] += 1
    else:
        metrics["wrong"] += 1
    
    # ALWAYS increment metrics["total"] for every sample.
    metrics["total"] += 1
    metrics["y_pred_list"].append(y_pred)

\end{minted}

This function is built to accept three positional arguments. Each argument plays a specific role in how the robot is drawn:
\begin{itemize}\setlength{\itemsep}{0pt}
    \item \texttt{i}: \\ An index or identifier for the robot.
    \item \texttt{canvas}: \\  The digital space where the robot will appear.
    \item \texttt{config}: \\ Configuration details for the robot's appearance.
    \item \texttt{colors}: \\  A collection of colors used for different parts of the robot.
    \item \texttt{body\_color}: \\  The specific color for the robot's body.
    \item \texttt{antenna\_color}: \\ The color for the robot's antenna.
\end{itemize}

\subsection{Best Practices: Enhancing Code Clarity with Descriptive Names and Docstrings}

In scenarios with multiple positional arguments (2 or more), the significance of clear \texttt{docstring} and clear naming variables becomes even more importance. These practices help prevent confusion and errors, ensuring that anyone using the function (e.g., update\_metrics), understands what each argument represents and in what order they should be provided.

Yet, it's important to recognize that mistakes can still occur. Users might inadvertently supply incorrect input or mix up the order of arguments. So, I hope this higlight the important ensuring supplying the right value and order for correct operation of the function's operation.

\subsection{Self Reflection: Your Experience with Positional Arguments}
Nevertheless, either you realise it or not, but you actually have been working with the concept of positional arguments since the third session of our tutorials. So Kudos to you for that. Therefore, I did not design any specific task for this section, as I believe you already have a good grasp of how positional arguments work in \texttt{Python} functions.






\newpage


\section{A Reusable Line-by-Line Debugging Workflow}
\begin{tcolorbox}[title=Line-by-Line Debugging Workflow]
When students are unsure how an argument moves through the program, use this debugging routine:
\begin{enumerate}
    \item Set a breakpoint on the line in \texttt{main()} that calls the function you want to study.
    \item Start the debugger and stop at that line before the function call happens.
    \item Read the argument values from left to right.
    \item Use \textbf{Step Into} to enter the function.
    \item Compare the values at the call site with the parameter names in the function header.
    \item Keep a close eye on variables such as \texttt{filepath}, \texttt{threshold}, \texttt{print\_each}, and \texttt{use\_fit}.
    \item If the behavior looks wrong, ask: ``Did the wrong value enter the wrong parameter?'' This question solves many beginner argument bugs.
\end{enumerate}
\end{tcolorbox}

\section{Summary and Review}
\begin{itemize}
    \item Positional arguments are matched by order.
    \item Keyword arguments are matched by parameter name.
    \item Default values let a function stay easy to call.
    \item Session 5 turns the classifier into a more configurable pipeline without changing the core prediction rule.
\end{itemize}

\textbf{Reflection Questions:}
\begin{enumerate}
    \item Why can positional arguments become harder to read when a function has many parameters?
    \item Why does \texttt{calculate\_accuracy(correct=..., total=...)} make the call site easier to understand?
    \item Why is the \texttt{use\_fit} flag a good example of controlling behavior from outside the function?
    \item When debugging, why is \textbf{Step Into} especially useful for studying arguments?
\end{enumerate}

\textbf{Micro-Exercise:} After students finish the tutorial, ask them to change the keyword override run from \texttt{threshold=2.2} to \texttt{threshold=1.5}. Then ask them to compare the predictions and summary with the fit-based run.
